La siguiente tabla sintética representa aproximaciones \(Q(h)\) a un mismo valor \(y(T)\), que se supone existente pero no se conoce exactamente. El instante final y el procedimiento se mantienen al cambiar \(h\). Las cuatro primeras salidas están redondeadas a seis decimales; las dos últimas, a tres.
\[
\begin{array}{c|rrrrrr}
h&1/2&1/4&1/8&1/16&1/32&1/64\\\hline
Q(h)&1.240000&1.210000&1.202500&1.200625&1.200&1.200
\end{array}
\]
A partir de las cuatro primeras aproximaciones, examine la hipótesis de trabajo \(Q(h)\approx L+Ch^p\), con \(C\ne0\) y \(p>0\). Estime \(p\), el límite \(L\) y el error absoluto de \(Q(1/16)\) bajo esa hipótesis. Se busca un error inferior a \(0.001\): indique qué conclusión está apoyada por el patrón y si los datos bastan para garantizarla. Explique qué información aportan las dos salidas finales iguales y qué comprobación adicional permitiría reforzar la conclusión sobre el error.
